%% ma: L12-B06-0074
%% lop: 12
%% chuong: 2
%% bai: 6
%% dang: 12.6.16
%% loai: TN4
%% muc: VD
%% dap_an: "C"
%% nguon: "CĐ2.1.pdf giáo viên gửi 27/10/2026 (M07, Lê Minh Khôi – Nguyễn Anh Khôi, Chương II Chủ đề 1 Vectơ trong không gian), Câu 34"
%% kien_thuc: "Bài 6 (tích vô hướng bằng 0 khi hai vectơ vuông góc), lớp 11 (đường thẳng vuông góc với mặt phẳng), lớp 10 (hệ thức lượng trong tam giác vuông)"
%% trang_thai: chinh-thuc
\begin{ex}
Cho tứ diện $ABCD$ có cạnh $AC\perp(BCD)$ và $BCD$ là tam giác đều cạnh bằng $1$. Biết $AC=\sqrt2$ và $M$ là trung điểm $BD$. Điểm $H$ thuộc $AM$ thỏa $\overrightarrow{CH}\cdot\overrightarrow{AM}=0$. Tính độ dài vectơ $\overrightarrow{CH}$.
\choice{$\dfrac{\sqrt6}{11}$}{$\dfrac4{11}$}{\True $\dfrac{\sqrt{66}}{11}$}{$\sqrt6$}
\loigiai{$CH\perp AM$ nên $CH$ là đường cao của tam giác $ACM$ vuông tại $C$, với $CM=\dfrac{\sqrt3}2$, $AM=\sqrt{2+\dfrac34}=\dfrac{\sqrt{11}}2$.

$CH=\dfrac{AC\cdot CM}{AM}=\dfrac{\sqrt2\cdot\frac{\sqrt3}2}{\frac{\sqrt{11}}2}=\dfrac{\sqrt{66}}{11}$.}
\end{ex}
